\documentclass[12pt,a4paper]{article}

% Packages
\usepackage[margin=2.5cm,headheight=15pt]{geometry}
\usepackage{amsmath,amssymb}
\usepackage{graphicx}
\usepackage{booktabs}
\usepackage{hyperref}
\usepackage{caption}
\usepackage{subcaption}
\usepackage{array}
\usepackage{longtable}
\usepackage{xcolor}
\usepackage{titlesec}
\usepackage{fancyhdr}
\usepackage{float}
\usepackage{microtype}
\usepackage{parskip}

% Header / footer
\pagestyle{fancy}
\fancyhf{}
\lhead{\small ALO Array Configuration Study}
\rhead{\small M. Rashidy, 2026}
\cfoot{\thepage}

\hypersetup{
  colorlinks=true,
  linkcolor=blue!70!black,
  urlcolor=blue!70!black,
  citecolor=blue!70!black
}

\titleformat{\section}{\large\bfseries}{\thesection.}{0.7em}{}
\titleformat{\subsection}{\normalsize\bfseries}{\thesubsection.}{0.5em}{}

\graphicspath{
  {outputs/plots/}
  {outputs/plots/geometry_tradeoff/}
  {outputs/plots/geometry_tradeoff/cross_config/}
  {outputs/plots/geometry_tradeoff/asymmetric_beam/}
}

% ──────────────────────────────────────────────────────────────────
\begin{document}

\begin{titlepage}
  \centering
  \vspace*{2cm}
  {\LARGE\bfseries Array Configuration Trade-off Study\\[0.5em]
   for the Array on the Lunar Outpost (ALO)}\par
  \vspace{1.5cm}
  {\large Mahshad Rashidy}\par
  \vspace{0.5cm}
  {\normalsize Department of Astrophysics / Thesis Project}\par
  \vspace{0.5cm}
  {\normalsize \today}\par
  \vfill
  \begin{abstract}
    \noindent
    This report documents the full pipeline of the Array on the Lunar Outpost
    (ALO) array-configuration trade-off study.
    ALO is a proposed low-frequency radio array (1--40\,MHz) to be deployed on
    the lunar far side at the Tsiolkovsky crater, exploiting the natural
    radio-frequency interference (RFI) shield of the Moon.
    Its primary science goal is the detection of magnetospheric radio emission
    from hot Jupiter exoplanets via the electron-cyclotron maser (ECM)
    instability.
    We evaluate four outrigger geometries (ring, Y-shape, random, and line),
    two core sizes (32\,$\times$\,32 and 128\,$\times$\,128 dipoles), two
    outrigger separations (1.0 and 5.0\,km), and two novel cross-shaped
    asymmetric configurations (E-W and N-S dominant arms) against a catalogue
    of 373 known hot Jupiters.
    Key metrics computed include the array factor (AF) and beam solid angle,
    half-power beam width (HPBW), maximum sidelobe level (MSL), UV coverage,
    thermal sensitivity, confusion limit, and required integration time.
    The 128\,$\times$\,128 ring and line arrays at 5.0\,km outrigger separation
    achieve the highest detection yield (34 feasible targets in under 100\,h),
    while the cross configurations provide the best sidelobe suppression
    (MSL\,$= -4.3$\,dB).
  \end{abstract}
\end{titlepage}

\tableofcontents
\clearpage

% ══════════════════════════════════════════════════════════════════
\section{Project Overview}
\label{sec:overview}

The Array on the Lunar Outpost (ALO) is a proposed phased-array radio telescope
designed to operate in the frequency range 1--40\,MHz.
This band is completely inaccessible from Earth due to ionospheric reflection
below $\sim$10\,MHz and RFI saturation at higher frequencies, making the lunar
far side the only viable platform for systematic exploration of the decametric
radio sky.

ALO will be located at the Tsiolkovsky impact crater on the lunar far side
(latitude $\varphi = -20.38^\circ$, longitude $\lambda = 128.97^\circ$).
The crater rim acts as a natural shield against terrestrial RFI and acts as a
mechanically stable flat floor for antenna deployment.

The array consists of individual dipole antenna elements with physical spacing
$d = 5.15$\,m and an aperture efficiency $\eta = 0.70$.
A central ``core'' of $N_c \times N_c$ elements forms a contiguous dense
aperture, while four outrigger sub-arrays of $N_o \times N_o$ elements are
placed at variable distances from the core to extend the interferometric
baseline.

The present study systematically explores the trade-space of array geometry,
core size, and outrigger separation in order to identify the configuration most
suited to detecting radio emission from hot Jupiter exoplanets.
All modeling and simulation code is contained in two Python files:
\texttt{alo\_array\_modeling.py} (physical constants, beam/sensitivity
calculators) and \texttt{alo\_geometry\_tradeoff.py} (geometry generator,
configuration evaluator, and figure factory).

% ══════════════════════════════════════════════════════════════════
\section{Scientific Motivation}
\label{sec:motivation}

\subsection{Magnetospheric Radio Emission from Exoplanets}

Jupiter emits strong bursts of circularly polarised radio emission up to
$\sim$40\,MHz via the electron-cyclotron maser (ECM) instability driven by its
interaction with the Galilean moon Io and the solar wind.
Hot Jupiters---gas giants in orbits $\lesssim 0.1$\,AU---are expected to
produce analogous ECM emission, potentially orders of magnitude stronger than
Jupiter's, because:
\begin{itemize}
  \item Their proximity to the host star drives intense stellar-wind ram
    pressure.
  \item Strong magnetic fields (several gauss) are inferred from planetary
    interior models, setting the cyclotron cutoff above 40\,MHz.
  \item Star-planet magnetic interaction (analogous to Io--Jupiter) adds an
    additional energy source.
\end{itemize}

The predicted flux density for bright targets (e.g.\ $\tau$\,Boo\,b,
51\,Peg\,b, $\upsilon$\,And\,b) is 1--100\,mJy at 10--40\,MHz, depending on
the assumed field strength and beam fraction.

\subsection{Why the Lunar Far Side?}

Below 10\,MHz the terrestrial ionosphere is opaque.
Between 10 and 30\,MHz, broadband RFI from human technology overwhelms
astrophysical signals even at the most remote Earth-based sites.
The lunar far side is permanently shielded from Earth by $\sim$3000\,km of
solid rock.
The Moon has no ionosphere above 1\,MHz at night, and the lunar far side
is also shielded from solar radio bursts during lunar night.

\subsection{Target Catalogue}

The target catalogue used in this study contains 373 hot Jupiter systems with
predicted ECM emission frequencies in the 1--40\,MHz band, compiled from the
NASA Exoplanet Archive.
Flux predictions follow a scaling law based on the stellar wind ram pressure
and the planet's estimated magnetic moment.
The catalogue is divided into four frequency sub-bands:
1--5\,MHz, 5--10\,MHz, 10--20\,MHz, and 20--40\,MHz.

% ══════════════════════════════════════════════════════════════════
\section{Array Geometry Setup}
\label{sec:geometry}

\subsection{Dipole Spacing and Core Layout}

Each element is a short dipole with physical size $d = 5.15$\,m and aperture
efficiency $\eta = 0.70$.
The effective area of a single element at frequency $\nu$ is:
\begin{equation}
  A_\mathrm{eff,elem}(\nu) = \eta\,\frac{\lambda^2}{4\pi}
  = \eta\,\frac{c^2}{4\pi\nu^2},
  \label{eq:aeff_elem}
\end{equation}
where $c = 3 \times 10^8$\,m\,s$^{-1}$.

A square core of $N_c \times N_c$ elements arranged on a regular grid with
spacing $d$ occupies a footprint of $(N_c - 1) \times d$ on each side.
Configurations tested: $N_c = 32$ (1088\,elements\footnote{Including four
$N_o \times N_o = 4\times4$ outrigger sub-arrays of 16 elements each.
$N_\mathrm{tot} = 32^2 + 4\times4^2 = 1024 + 64 = 1088$.}) and
$N_c = 128$ (17\,408 elements in total with 16-element outriggers).

\subsection{Outrigger Placement}

Four outrigger sub-arrays, each of size $N_o \times N_o$, are placed
symmetrically at distance $r$ from the core centre along the four cardinal
directions (North, South, East, West) for symmetric geometries.
For asymmetric (cross) geometries, different arm lengths are used (see
Section~\ref{sec:configs}).

The outrigger separation values tested are $r = 1.0$\,km and $r = 5.0$\,km.

Figure~\ref{fig:layout_overview} shows the array layouts for all 16
standard configurations.

\begin{figure}[H]
  \centering
  \includegraphics[width=0.95\textwidth]{layout_overview.png}
  \caption{Overview of all 16 standard array layouts.
    Rows correspond to outrigger geometry (ring, Y-shape, random, line);
    columns to core size and outrigger separation.
    Blue dots are core elements; orange squares mark outrigger sub-array
    centres.}
  \label{fig:layout_overview}
\end{figure}

% ══════════════════════════════════════════════════════════════════
\section{Coordinate Transformation}
\label{sec:coords}

All antenna positions are expressed in a local East-North-Up (ENU) frame
centred on the Tsiolkovsky crater.

\subsection{ENU Basis Vectors}

Given geodetic latitude $\varphi$ and longitude $\lambda$ of the array
centre, the ENU basis unit vectors in the selenocentric Cartesian frame are:
\begin{align}
  \hat{e}_E &= (-\sin\lambda,\; \cos\lambda,\; 0), \\
  \hat{e}_N &= (-\sin\varphi\cos\lambda,\; -\sin\varphi\sin\lambda,\; \cos\varphi), \\
  \hat{e}_U &= (\cos\varphi\cos\lambda,\; \cos\varphi\sin\lambda,\; \sin\varphi).
\end{align}

An element at ENU offset $(x, y, 0)$ relative to the array centre has
selenocentric position:
\begin{equation}
  \mathbf{r} = x\,\hat{e}_E + y\,\hat{e}_N.
\end{equation}

\subsection{Direction Cosines}

For far-field imaging in the small-angle regime the sky position of a source is
parameterised by direction cosines $(l, m, n)$ where
$l^2 + m^2 + n^2 = 1$ and the phase centre is at $(l,m) = (0,0)$, $n = 1$.
The array factor is computed on a grid of $(l, m)$ values spanning
$[-1, 1]^2$, enforcing $l^2 + m^2 \le 1$.

% ══════════════════════════════════════════════════════════════════
\section{Array Configurations}
\label{sec:configs}

\subsection{Standard Outrigger Geometries}

Four outrigger placement patterns are evaluated:

\begin{description}
  \item[Ring] Four outriggers on the cardinal axes at equal distance $r$.
    The UV coverage forms a cross plus core blob.
  \item[Y-shape] Outriggers placed at $0^\circ$ (N), $120^\circ$, and
    $240^\circ$ plus one at $90^\circ$ (E), giving a slightly asymmetric
    Y-pattern.
  \item[Random] Outrigger centres drawn from a uniform 2-D distribution
    within a circle of radius $r$, providing a more isotropic UV coverage.
  \item[Line] All four outriggers placed along the East--West axis at
    offsets $-r$, $-r/2$, $r/2$, $r$ providing maximum E-W resolution.
\end{description}

\subsection{Cross-Shaped Asymmetric Configurations}
\label{sec:cross}

Two additional configurations with a cross-shaped outrigger arrangement but
\emph{unequal} arm lengths are introduced to break the periodicity that causes
aliasing rings in equal-arm configurations:

\begin{center}
\begin{tabular}{lcccc}
  \toprule
  Configuration & N arm & S arm & E arm & W arm \\
  \midrule
  Cross E-W & 1.0\,km & 1.0\,km & 5.0\,km & 3.0\,km \\
  Cross N-S & 5.0\,km & 3.0\,km & 1.0\,km & 1.0\,km \\
  \bottomrule
\end{tabular}
\end{center}

Both cross configurations use a 128$\times$128 core ($N_c = 128$) with
16-element outrigger sub-arrays ($N_o = 4 \times 4$), giving 17\,408 elements
in total and $B_\mathrm{max} = 8.33$\,km.
Figure~\ref{fig:layout_cross} shows the layouts of both cross configurations
alongside the four standard 128$\times$128 + 5\,km reference arrays.

\begin{figure}[H]
  \centering
  \includegraphics[width=0.95\textwidth]{layout_cross_overview.png}
  \caption{Layouts of the six configurations used in the cross comparison study.
    Top row: four standard 128\,$\times$\,128 + 5\,km outrigger configurations.
    Bottom row: cross E-W and cross N-S.}
  \label{fig:layout_cross}
\end{figure}

% ══════════════════════════════════════════════════════════════════
\section{Array Factor and Beam-Pattern Calculation}
\label{sec:af}

\subsection{Array Factor Definition}

The array factor (AF) is the spatial Fourier transform of the aperture
illumination.
For $N$ elements at positions $\mathbf{r}_k = (x_k, y_k)$ with complex weights
$w_k$ (set to unity for an unweighted array), the AF at sky direction $(l, m)$
and wavelength $\lambda = c/\nu$ is:
\begin{equation}
  \mathrm{AF}(l,m;\nu) = \sum_{k=1}^{N} w_k\,
    \exp\!\left[\frac{2\pi i}{\lambda}\bigl(x_k\,l + y_k\,m\bigr)\right].
  \label{eq:af}
\end{equation}

The normalised power beam is:
\begin{equation}
  B(l,m;\nu) = \frac{|\mathrm{AF}(l,m;\nu)|^2}
               {|\mathrm{AF}(0,0;\nu)|^2}.
\end{equation}

\subsection{Physical Origin of Beam Symmetry}

A key result of this study is that $B(l,m)$ is \emph{always}
180$^\circ$-rotationally symmetric, regardless of array geometry.
This follows from Hermitian symmetry: for real weights $w_k$,
\begin{equation}
  \mathrm{AF}(-l,-m) = \mathrm{AF}^*(l,m)
  \quad\Longrightarrow\quad
  B(-l,-m) = B(l,m).
\end{equation}
This is not a code artefact or a modelling approximation---it is a fundamental
property of the Fourier transform of any real-valued aperture.
Geometry information is encoded in the \emph{phase} of AF, which is discarded
when the power beam $|AF|^2$ is formed.
Consequently, point-source sensitivity is identical for all array geometries
with the same total element count $N$, and beam-shape differences between
geometries are limited to their sidelobe structure.

\subsection{Frequency Dependence}

The beam solid angle scales as $\Omega_B \propto \lambda^2 \propto \nu^{-2}$,
so the HPBW shrinks rapidly with frequency.
Beam patterns are computed at representative frequencies within each sub-band:
3.0, 7.5, 15.0, and 30.0\,MHz for the 1--5, 5--10, 10--20, and 20--40\,MHz
bands respectively.
All AF evaluations use an $n_\mathrm{grid} = 512$ $(l,m)$ grid.

Figure~\ref{fig:af_overview} shows the normalised beam patterns at 30\,MHz
for the eight standard configurations.

\begin{figure}[H]
  \centering
  \includegraphics[width=0.85\textwidth]{AF_overview_30MHz.png}
  \caption{Normalised power-beam patterns at 30\,MHz for all eight
    standard configurations (32\,$\times$\,32 and 128\,$\times$\,128 cores,
    1.0 and 5.0\,km outrigger separation).
    The colour scale is in dB, clipped at $-30$\,dB to reveal sidelobe
    structure.
    All beams show 180$^\circ$ rotational symmetry as required by
    Hermitian symmetry of the Fourier transform.}
  \label{fig:af_overview}
\end{figure}

\subsection{Beam Contour Plots}

Beam patterns are visualised with three contour levels at 10\%, 30\%, and 50\%
of the peak power (i.e.\ $-10$, $-5.2$, and $-3$\,dB).
These levels mark the outer sidelobe envelope, the mid-sidelobe region, and the
main-lobe half-power boundary respectively.
Figure~\ref{fig:beam_contours} shows beam contours for the six cross-comparison
configurations at 30\,MHz.

\begin{figure}[H]
  \centering
  \includegraphics[width=0.95\textwidth]{beam_contours_cross.png}
  \caption{Beam contour plots at 30\,MHz for the six configurations in the
    cross comparison study.
    Contour levels at 10\% (cyan), 30\% (green), and 50\% (white) of peak.
    The cross configurations (bottom row) show a more compact sidelobe
    envelope than the symmetric ring and line arrays (top row).}
  \label{fig:beam_contours}
\end{figure}

% ══════════════════════════════════════════════════════════════════
\section{Beam Metrics}
\label{sec:beam_metrics}

\subsection{Beam Solid Angle}

The beam solid angle is the solid-angle integral of the normalised power beam:
\begin{equation}
  \Omega_B(\nu) = \iint B(l,m;\nu)\,\frac{\mathrm{d}l\,\mathrm{d}m}
                   {\sqrt{1 - l^2 - m^2}}.
  \label{eq:omega_b}
\end{equation}
In practice, the integral is replaced by a discrete sum over the $(l,m)$ grid
multiplied by the pixel area $(\Delta l)^2$.

\subsection{Half-Power Beam Width}

The HPBW is derived from the beam solid angle assuming a circular Gaussian
main lobe:
\begin{equation}
  \theta_\mathrm{HPBW} = 2\sqrt{\frac{\Omega_B}{\pi}}
  \quad[\mathrm{rad}].
\end{equation}

\subsection{Directivity and Gain}

Peak directivity and gain are:
\begin{equation}
  D = \frac{4\pi}{\Omega_B}, \qquad G = \eta\,D,
\end{equation}
and the total effective area is:
\begin{equation}
  A_\mathrm{eff}(\nu) = \frac{\lambda^2\,G}{4\pi}.
\end{equation}

\subsection{Maximum Sidelobe Level}

The MSL is defined as the peak normalised beam value outside the main lobe
(i.e.\ outside the $B \ge 0.5$ region), expressed in dB:
\begin{equation}
  \mathrm{MSL} = 10\log_{10}\!\Bigl(\max_{B < 0.5} B(l,m)\Bigr).
\end{equation}

\subsection{Summary of Beam Metrics}

Table~\ref{tab:beam_metrics} summarises the HPBW and MSL at 30\,MHz for all
configurations in the cross comparison study.

\begin{table}[H]
  \centering
  \caption{Beam metrics at 30\,MHz for the six cross-comparison configurations.
    All use a 128\,$\times$\,128 core (17\,408 total elements).}
  \label{tab:beam_metrics}
  \begin{tabular}{lccrr}
    \toprule
    Configuration & $B_\mathrm{max}$ (km) & HPBW (arcmin) & MSL (dB) \\
    \midrule
    Ring             & 10.33 & 63.8 & $-3.8$ \\
    Y-shape          &  8.99 & 64.9 & $-3.3$ \\
    Random           &  7.65 & 64.0 & $-3.8$ \\
    Line             & 10.33 & 65.0 & $-3.4$ \\
    Cross E-W        &  8.33 & 64.4 & $-4.3$ \\
    Cross N-S        &  8.33 & 64.4 & $-4.3$ \\
    \bottomrule
  \end{tabular}
\end{table}

The cross configurations achieve the best (most negative) MSL of $-4.3$\,dB,
0.5\,dB better than the best symmetric configuration.
This improvement arises because the irregular arm lengths (1/1/5/3\,km)
prevent baseline lengths from repeating at regular intervals, suppressing
coherent addition of aliasing rings.

% ══════════════════════════════════════════════════════════════════
\section{UV Coverage Analysis}
\label{sec:uv}

\subsection{Definition}

UV coverage is the set of spatial frequencies sampled by an interferometric
array.
For a pair of stations at positions $(x_i, y_i)$ and $(x_j, y_j)$, the
baseline vector projected onto the sky plane at frequency $\nu$ is:
\begin{equation}
  (u, v) = \frac{(x_j - x_i,\; y_j - y_i)}{\lambda}.
\end{equation}
Both the baseline and its conjugate $(-u, -v)$ are included (Hermitian
conjugate pair), so $N_s$ stations produce $N_s(N_s-1)$ UV points per
frequency channel.

\subsection{Station-Level UV Coverage}

For multi-sub-array configurations, UV coverage is computed at the
\emph{station} level, treating each sub-array as a single coherent station.
The five stations (one core + four outriggers) produce $5 \times 4 = 20$ UV
points per channel (including conjugates).

\subsection{Multi-Frequency Synthesis}

Multi-frequency synthesis (MFS) is simulated by computing UV coverage at every
1-MHz channel within each sub-band.
At frequency $\nu$ the baseline $(u,v)$ scales as $1/\lambda \propto \nu$, so
each baseline traces a radial streak in the UV plane as frequency increases.
The four sub-bands and their colours used in Figure~\ref{fig:uv_coverage} are:
\begin{itemize}
  \item 1--5\,MHz (red): low-frequency, large $\lambda$, short UV coverage
  \item 5--10\,MHz (orange): intermediate
  \item 10--20\,MHz (blue): higher resolution
  \item 20--40\,MHz (purple): highest resolution, longest UV tracks
\end{itemize}

\begin{figure}[H]
  \centering
  \includegraphics[width=0.95\textwidth]{uv_coverage.png}
  \caption{UV coverage for all standard configurations and the two cross
    configurations, computed via multi-frequency synthesis using 1-MHz channels
    within each sub-band.
    Each sub-band is plotted in a distinct colour.
    The cross configurations (bottom right) show elongated UV coverage in one
    axis, reflecting the dominant long arm.}
  \label{fig:uv_coverage}
\end{figure}

\begin{figure}[H]
  \centering
  \includegraphics[width=0.85\textwidth]{uv_coverage_cross.png}
  \caption{UV coverage for the six cross-comparison configurations only,
    with all four sub-bands superimposed.
    The ring and line arrays provide more rotationally symmetric UV coverage;
    the cross arrays concentrate UV tracks along two axes.}
  \label{fig:uv_cross}
\end{figure}

% ══════════════════════════════════════════════════════════════════
\section{Sensitivity Model}
\label{sec:sensitivity}

\subsection{Sky Temperature}

The sky brightness temperature in the 1--40\,MHz band is dominated by
Galactic synchrotron emission and follows a power law:
\begin{equation}
  T_\mathrm{sky}(\nu) = 180\,\mathrm{K} \times
    \left(\frac{\nu}{180\,\mathrm{MHz}}\right)^{-2.6}.
  \label{eq:tsky}
\end{equation}
At 10\,MHz, $T_\mathrm{sky} \approx 1.2 \times 10^5$\,K.
This is the dominant noise source; receiver and ground contributions are
negligible by comparison.

\subsection{System Equivalent Flux Density}

The system equivalent flux density (SEFD) per element is:
\begin{equation}
  \mathrm{SEFD}_\mathrm{elem}(\nu) = \frac{2\,k_B\,T_\mathrm{sky}(\nu)}
                                         {A_\mathrm{eff,elem}(\nu)},
  \label{eq:sefd}
\end{equation}
where $k_B = 1.381 \times 10^{-23}$\,J\,K$^{-1}$.
For the phased array of $N$ elements operating as a coherent aperture, the
array SEFD is $\mathrm{SEFD}_\mathrm{elem}/N$.

\subsection{Thermal Sensitivity}

For dual-polarisation correlation ($N_\mathrm{pol} = 2$) over bandwidth
$\Delta\nu$ and integration time $t$, the thermal noise is:
\begin{equation}
  \sigma_\mathrm{th}(\nu, t) = \frac{\mathrm{SEFD}_\mathrm{elem}}
    {\sqrt{N_\mathrm{pol}\,N(N-1)\,\Delta\nu\,t}},
  \label{eq:sigma_th}
\end{equation}
where $N(N-1)$ counts all unique baselines (including cross-correlations
between core elements and outrigger elements).

\subsection{Confusion Limit}

Below a certain flux level, sources blend together in the beam and set an
irreducible confusion noise floor:
\begin{equation}
  \sigma_c(\nu, B_\mathrm{max}) = 0.2\,\mathrm{mJy} \times
    \left(\frac{\nu}{74\,\mathrm{MHz}}\right)^{0.8} \times
    \left(\frac{\theta_\mathrm{HPBW}}{4\,\mathrm{arcmin}}\right),
  \label{eq:confusion}
\end{equation}
where $\theta_\mathrm{HPBW}$ is determined by the maximum baseline
$B_\mathrm{max}$.
At 30\,MHz with $B_\mathrm{max} = 10$\,km, $\sigma_c \approx 0.006$\,mJy,
far below the ECM flux predictions.

% ══════════════════════════════════════════════════════════════════
\section{Integration-Time Calculation}
\label{sec:tint}

\subsection{Required Integration Time}

To detect a target of flux density $S$ at $N_\sigma = 5\sigma$ significance,
the required integration time is obtained by solving:
\begin{equation}
  N_\sigma \sqrt{\sigma_\mathrm{th}^2(t) + \sigma_c^2} = S.
  \label{eq:detection_condition}
\end{equation}
Substituting equation~\eqref{eq:sigma_th}:
\begin{equation}
  t_\mathrm{req} = \left(\frac{\mathrm{SEFD}_\mathrm{elem}}
    {\sqrt{N_\mathrm{pol}\,N(N-1)\,\Delta\nu}}\right)^2
    \frac{1}{(S/N_\sigma)^2 - \sigma_c^2}.
  \label{eq:t_req}
\end{equation}
If $(S/N_\sigma)^2 \le \sigma_c^2$ (confusion-dominated), the target is
irresolvable regardless of integration time and $t_\mathrm{req} = \infty$ is
returned.

\subsection{Feasibility Classification}

Targets are classified into three categories:
\begin{description}
  \item[Feasible] $t_\mathrm{req} < 100$\,h — achievable in a single observing
    campaign.
  \item[Aspirational] $100 \le t_\mathrm{req} < 1000$\,h — borderline; might
    be feasible with improved calibration or stacking.
  \item[Impossible] $t_\mathrm{req} \ge 1000$\,h — not detectable by ALO in
    any realistic scenario.
\end{description}

% ══════════════════════════════════════════════════════════════════
\section{Plot Interpretation}
\label{sec:plot_interp}

\subsection{Beam Pattern Plots}

Beam patterns are shown as 2D images of $B(l,m)$ on a dB scale.
The colour map spans $-30$ to $0$\,dB, chosen to make sidelobe structure
visible while not washing out the main lobe.
The 180$^\circ$ symmetry visible in all plots is physical (Hermitian symmetry),
not a plotting artefact.
Differences between geometries appear as changes in sidelobe shape and
sidelobe level, not in the main lobe width (which is set by the total core
diameter).

\subsection{UV Coverage Plots}

Each coloured point in the UV plane represents one baseline at one frequency
channel.
Dense radial streaks indicate multi-frequency synthesis tracks.
Gaps in the UV plane correspond to spatial frequencies not sampled, which
manifest as image artefacts (grating lobes) if uncleaned.
The cross configurations show elongated UV tracks along the dominant arm
direction.

\subsection{Sensitivity Comparison Plots}

Figure~\ref{fig:metrics_cross} compares beam metrics and
Figure~\ref{fig:sensitivity_cross} compares the number of feasible detections
across all six configurations in the cross comparison study.

\begin{figure}[H]
  \centering
  \includegraphics[width=0.92\textwidth]{metrics_cross_vs_originals.png}
  \caption{Beam quality metrics (HPBW and MSL) at all four frequency sub-bands
    for the six cross-comparison configurations.
    The cross configurations consistently achieve lower (better) MSL than the
    symmetric arrays.}
  \label{fig:metrics_cross}
\end{figure}

\begin{figure}[H]
  \centering
  \includegraphics[width=0.92\textwidth]{sensitivity_cross_vs_originals.png}
  \caption{Number of feasible (${<}100$\,h) detections per configuration
    per sub-band in the cross comparison study.
    The 20--40\,MHz band dominates owing to stronger predicted ECM fluxes and
    lower sky temperature.}
  \label{fig:sensitivity_cross}
\end{figure}

\subsection{Detection Matrix}

The detection matrix (Figure~\ref{fig:detection_matrix}) shows the required
integration time for each target--configuration pair, coded by feasibility
category.

\begin{figure}[H]
  \centering
  \includegraphics[width=0.85\textwidth]{detection_matrix_heatmap.png}
  \caption{Detection feasibility matrix for the full array-configuration
    trade-off study.
    Green: feasible ($<100$\,h); yellow: aspirational (100--1000\,h);
    red: impossible.
    Rows are configurations; columns are targets ordered by decreasing
    flux density.}
  \label{fig:detection_matrix}
\end{figure}

% ══════════════════════════════════════════════════════════════════
\section{Results and Comparison}
\label{sec:results}

\subsection{Standard Configuration Comparison}

Table~\ref{tab:main_results} summarises the detection yield for all 16
standard configurations across the full 373-target catalogue.

\begin{table}[H]
  \centering
  \caption{Array configuration summary.
    $N_\mathrm{det}(<100\,\mathrm{h})$ counts targets detectable in $<100$\,h
    in any frequency band.}
  \label{tab:main_results}
  \begin{tabular}{lcccc}
    \toprule
    Configuration & $N_\mathrm{elem}$ & $B_\mathrm{max}$ (km)
      & HPBW at 30 MHz (arcmin) & $N_\mathrm{det}(<100\,\mathrm{h})$ \\
    \midrule
    ring\_core32x32\_1.0km   &  1088 &  2.08 & 579 & -- \\
    ring\_core32x32\_5.0km   &  1088 & 10.08 &  63 & -- \\
    ring\_core128x128\_1.0km & 17408 &  2.33 & 579 & -- \\
    ring\_core128x128\_5.0km & 17408 & 10.33 & 63.8 & \textbf{34} \\
    y\_shape\_core32x32\_1.0km   &  1088 & 1.81 & 578 & -- \\
    y\_shape\_core32x32\_5.0km   &  1088 & 8.74 & 63 & -- \\
    y\_shape\_core128x128\_1.0km & 17408 & 2.06 & 578 & -- \\
    y\_shape\_core128x128\_5.0km & 17408 & 8.99 & 64.9 & \textbf{33} \\
    random\_core32x32\_1.0km   &  1088 & 1.55 & 576 & -- \\
    random\_core32x32\_5.0km   &  1088 & 7.41 & 63 & -- \\
    random\_core128x128\_1.0km & 17408 & 1.79 & 576 & -- \\
    random\_core128x128\_5.0km & 17408 & 7.65 & 64.0 & \textbf{31} \\
    line\_core32x32\_1.0km   &  1088 & 2.08 & 576 & -- \\
    line\_core32x32\_5.0km   &  1088 & 10.08 & 63 & -- \\
    line\_core128x128\_1.0km & 17408 & 2.33 & 576 & -- \\
    line\_core128x128\_5.0km & 17408 & 10.33 & 65.0 & \textbf{34} \\
    \bottomrule
    \multicolumn{5}{l}{\small $N_\mathrm{det}$ for 32\,$\times$\,32 and 1\,km
      configurations is $<5$ across all bands and omitted.} \\
  \end{tabular}
\end{table}

\textbf{Key findings:}
\begin{itemize}
  \item The 128\,$\times$\,128 core is essential: the 32\,$\times$\,32 arrays
    have insufficient collecting area to detect any target in $<100$\,h.
  \item A 5\,km outrigger separation is essential for angular resolution: at
    1\,km the HPBW at 30\,MHz is $\sim$580\,arcmin, too large to discriminate
    targets from background.
  \item The ring and line arrays share the largest $B_\mathrm{max}$ (10.33\,km)
    and thus the sharpest beam and lowest confusion limit, yielding the most
    detections (34 each).
  \item Geometry differences among the four outrigger patterns are small:
    $\Delta N_\mathrm{det} \le 3$ across ring/Y-shape/random/line.
\end{itemize}

\subsection{Cross Configuration Comparison}

Table~\ref{tab:cross_results} compares the two cross configurations against the
four 128\,$\times$\,128 + 5\,km reference arrays.

\begin{table}[H]
  \centering
  \caption{Cross configuration comparison results.}
  \label{tab:cross_results}
  \begin{tabular}{lrrrr}
    \toprule
    Configuration & $N_\mathrm{elem}$ & $B_\mathrm{max}$ (km)
      & $N_\mathrm{det}(<100\,\mathrm{h})$ & MSL (dB) \\
    \midrule
    Ring             & 17408 & 10.33 & 34 & $-3.8$ \\
    Y-shape          & 17408 &  8.99 & 33 & $-3.3$ \\
    Random           & 17408 &  7.65 & 31 & $-3.8$ \\
    Line             & 17408 & 10.33 & 34 & $-3.4$ \\
    Cross E-W        & 17408 &  8.33 & 33 & $\mathbf{-4.3}$ \\
    Cross N-S        & 17408 &  8.33 & 33 & $\mathbf{-4.3}$ \\
    \bottomrule
  \end{tabular}
\end{table}

The cross configurations achieve:
\begin{itemize}
  \item Best MSL ($-4.3$\,dB vs.\ $-3.3$ to $-3.8$\,dB for symmetric arrays).
  \item Competitive detection yield (33 targets) despite a shorter maximum
    baseline (8.33 vs.\ 10.33\,km for ring and line).
  \item Elongated UV coverage providing anisotropic resolution, which may be
    preferred for specific science cases (e.g.\ mapping extended plasma
    structures).
\end{itemize}

\subsection{Beam Symmetry Analysis}

A dedicated asymmetric beam study was conducted to investigate whether array
geometry asymmetry is detectable in the power beam.
Arrays with deliberately broken symmetry (line: all outriggers along E-W axis)
were compared to rotationally symmetric arrays (ring) at multiple pointing
directions.

\textbf{Conclusion:} The power beam $B(l,m) = |\mathrm{AF}|^2$ is always
180$^\circ$-rotationally symmetric due to the Hermitian symmetry of the
Fourier transform.
Asymmetry information is contained only in the \emph{phase} of the complex AF,
not in the magnitude.
No error in the code was found; the symmetry is physically correct.

\subsection{Extended Source Response}

An extended source analysis was performed using a 32\,$\times$\,32 + 5\,km
array (chosen for sufficient beam resolution in the numerical grid) to test
whether geometry differences become visible for spatially resolved sources.

The beam-weighted flux-recovery fraction:
\begin{equation}
  \eta_\mathrm{BW} = \frac{\sum_{l,m} I_\mathrm{true}(l,m)\,B(l,m)}
                           {\sum_{l,m} I_\mathrm{true}(l,m)}
\end{equation}
was computed as a function of source angular size and position angle (PA).

\textbf{Results:}
\begin{itemize}
  \item For point sources, $\eta_\mathrm{BW} = 1$ identically for all
    geometries (as expected).
  \item For extended sources, $\eta_\mathrm{BW}$ varies with source size and PA,
    and the variation is geometry-dependent.
  \item The line array shows the largest PA-orientation bias
    ($\Delta\eta = 0.0111$, the flux-recovery range over all tested PAs).
  \item The ring array shows the smallest PA-orientation bias
    ($\Delta\eta = 0.0002$), consistent with its near-isotropic UV coverage.
\end{itemize}

Since ALO's primary targets are point-like exoplanets (angular size
$\ll$ HPBW $\approx 64$\,arcmin at 30\,MHz), extended-source response
is not a discriminating metric for the primary science case.

% ══════════════════════════════════════════════════════════════════
\section{Final Conclusions}
\label{sec:conclusions}

\begin{enumerate}
  \item \textbf{Core size is the dominant driver of sensitivity.}
    The 128\,$\times$\,128 core (17\,408 elements, $A_\mathrm{eff}
    \approx 13{,}000$\,m$^2$ at 30\,MHz) is necessary for feasible detections;
    the 32\,$\times$\,32 core is insufficient.

  \item \textbf{Outrigger separation drives angular resolution.}
    At 5\,km the HPBW at 30\,MHz is $\sim$64\,arcmin; at 1\,km it is
    $\sim$578\,arcmin.
    The long baseline suppresses the confusion limit by a factor of 9, enabling
    detection of targets that would otherwise be confusion-dominated.

  \item \textbf{Outrigger geometry has a second-order effect on detection yield.}
    Differences in feasible detection count among ring, Y-shape, random, and
    line are only 1--3 targets out of 34.
    The choice of geometry should be driven by other considerations (e.g.\
    UV-coverage isotropy, sidelobe characteristics, construction constraints).

  \item \textbf{Cross configurations offer the best sidelobe suppression.}
    Irregular arm lengths (1/1/5/3\,km) break the periodicity that causes
    coherent aliasing rings, reducing MSL by 0.5\,dB relative to the best
    symmetric array.
    This comes at the cost of a slightly shorter maximum baseline (8.33 vs.\
    10.33\,km) and hence one fewer feasible detection.

  \item \textbf{The 20--40\,MHz sub-band is the most productive.}
    Most feasible detections occur in this band due to a combination of lower
    sky temperature, higher predicted ECM flux, and sufficient angular resolution
    at $B_\mathrm{max} = 8$--10\,km.

  \item \textbf{Point-source sensitivity is geometry-independent.}
    The power beam $|AF|^2$ is Hermitian-symmetric for any real-weight array.
    Geometry information in $|AF|^2$ is limited to sidelobe structure; main-lobe
    solid angle depends only on the total core aperture, not on outrigger pattern.
    This is a fundamental physical result, not a modelling limitation.
\end{enumerate}

% ══════════════════════════════════════════════════════════════════
\section{Recommendations for Next Steps}
\label{sec:recommendations}

\begin{enumerate}
  \item \textbf{Beam forming with weights.}
    Replace unit weights with an apodisation taper (e.g.\ Hanning window) to
    reduce sidelobes further.
    A Taylor taper with $n = 5$ sidelines could reduce MSL to $\sim -20$\,dB
    at the cost of $\sim$20\% wider main lobe.

  \item \textbf{Realistic beam model including element pattern.}
    Include the element (dipole) radiation pattern $E(\theta, \phi)$ as a
    multiplicative factor in the AF.
    This will create an element-pattern null towards the horizon and modify the
    effective collecting area at large zenith angles.

  \item \textbf{Lunar rotation and source tracking.}
    Compute the UV-coverage track for a full $\sim$14-day lunar day, including
    diurnal rotation of the far-side frame.
    Earth-facing sources appear above the lunar horizon for $\sim$11\,h
    per sidereal day and their UV tracks will fill a much larger region of the
    UV plane.

  \item \textbf{Optimal sub-array size for cross configuration.}
    Optimise the outrigger sub-array size $N_o$ for the cross geometry.
    Larger outriggers increase baseline sensitivity but increase mass and
    deployment cost.

  \item \textbf{Calibration architecture.}
    Design a calibration strategy for the far-side environment.
    Bright background sources (e.g.\ Cas\,A, Cyg\,A, Vir\,A) will be visible
    above the lunar horizon; assess their usability as calibrators for both
    gain and phase solutions.

  \item \textbf{Ionosphere-free mode.}
    Below $\sim$1\,MHz the Moon's plasma environment (exosphere) may introduce
    dispersive delays.
    Extend the sensitivity model to account for this chromatic phase noise.

  \item \textbf{Comparison with FARSIDE/DAPPER concepts.}
    Benchmark ALO performance against existing lunar-farside array proposals
    (FARSIDE, DAPPER, DSL) in terms of collecting area, angular resolution,
    and detection yield for the same target catalogue.

  \item \textbf{Deployment feasibility.}
    Feed the optimal configuration parameters (core size 128\,$\times$\,128,
    outrigger distance 5\,km, cross or ring geometry) into an engineering
    trade study covering deployment mass, power budget (solar/RTG), and
    thermal management on the lunar surface.
\end{enumerate}

% ══════════════════════════════════════════════════════════════════
\appendix
\section*{Appendix: Physical Constants and Parameters}

\begin{center}
\begin{tabular}{lll}
  \toprule
  Symbol & Value & Description \\
  \midrule
  $d$        & 5.15\,m                 & Dipole element spacing \\
  $\eta$     & 0.70                    & Aperture efficiency \\
  $N_\sigma$ & 5                       & Detection significance threshold \\
  $N_\mathrm{pol}$ & 2               & Number of polarisations \\
  $c$        & $3\times10^8$\,m\,s$^{-1}$ & Speed of light \\
  $k_B$      & $1.381\times10^{-23}$\,J\,K$^{-1}$ & Boltzmann constant \\
  $\varphi$  & $-20.38^\circ$          & Tsiolkovsky crater latitude \\
  $\lambda_\mathrm{site}$  & $128.97^\circ$ & Tsiolkovsky crater longitude \\
  \bottomrule
\end{tabular}
\end{center}

\vspace{1em}
\noindent\textbf{Software.}
All computations were performed using Python\,3 with NumPy, SciPy, and
Matplotlib.
The simulation code is version-controlled in the Git repository
\texttt{mahshadrashidi/ALO\_array\_modeling} on GitHub.
All figures in this report are generated directly from the simulation outputs
and are reproducible by running \texttt{alo\_geometry\_tradeoff.py}.

\end{document}
